\section{Experiments, Demonstrations, and Results}
\label{sec:experiments}

\subsection{Evaluation questions}

The current evidence is organized around five evaluation questions.

\begin{description}[style=nextline]
  \item[EQ1: Package correctness]
  Do strict load, validation, mutation, persistence, and import operations reject
  malformed or unsafe data while preserving valid semantics?

  \item[EQ2: Semantic runtime correctness]
  Do docking, undocking, assemblies, events, metrics, and graph views remain mutually
  consistent under success and failure?

  \item[EQ3: Backend realization]
  Does the same semantic contract execute through the mock and \mujoco{} boundaries,
  including physical resource refusal and runtime constraint release?

  \item[EQ4: Physical actuation]
  Can articulated wheel effort, gravity, tire contact, measured connector frames, and
  runtime welds complete a dock/undock cycle without root motion injection?

  \item[EQ5: Multi-module reconfiguration]
  Can the physics path execute all four Driver-to-Snake connector replacements,
  including motion of connected three-module components, while preserving expected
  topology and bounded state?
\end{description}

The implementation does not yet answer statistical robustness, hardware fidelity, or
large-scale performance questions. Those are explicitly deferred to
\cref{sec:future}.

\subsection{Verification environment and reproducibility}

The current report snapshot is \modsim{} 0.1.0 at commit \code{4a5ca00}. A full
repository test run under Python 3.12 with the Studio and \mujoco{} dependencies
installed collected and passed 537 tests in 166.91 seconds. Branch-aware coverage of
the backend-neutral \code{modsim} package was 86\%. The wall time is reported only as
the cost of this verification run; it is not a controlled physics benchmark because
machine load, GUI libraries, coverage instrumentation, and test composition affect it.

Both committed example packs passed fresh simulation-profile structural validation
with zero errors, warnings, or informational issues:

\begin{lstlisting}[language=bash]
modsim pack validate examples/robot_packs/generic_cube --profile simulation
modsim pack validate examples/robot_packs/smores_ep --profile simulation
\end{lstlisting}

The test suite includes core unit tests, loader and writer fault cases, schema and
validator cases, mock/\mujoco{} conformance, Studio project and logging tests,
Runtime Inspector protocol/process/widget tests, and physical end-to-end tests. CI
separates core Python 3.11/3.12, headless \mujoco{}, Xvfb Studio, and packaging jobs.

\subsection{Experiment matrix}

\begin{table}[htbp]
\centering
\caption{Named demonstration matrix. Durations and steps are simulated quantities.}
\label{tab:demo-matrix}
\begin{tabularx}{\textwidth}{L{0.24\textwidth}L{0.19\textwidth}L{0.10\textwidth}L{0.18\textwidth}Y}
\toprule
Demo & Motion model & Modules & Duration / step & Expected edge history \\
\midrule
\code{dock} & kinematic approach & 2 & 4 s / 0.002 s & $0\rightarrow1$ \\
\code{dock\_undock} & kinematic approach/retract & 2 & 6 s / 0.002 s & $0\rightarrow1\rightarrow0$ \\
\code{smores\_diff\_drive\_dock\_undock} & wheel/contact physics & 2 & 12 s / 0.002 s & $0\rightarrow1\rightarrow0$ \\
\code{smores\_driver\_to\_snake} & kinematic component staging & 7 & 14 s / 0.002 s & four $6\rightarrow5\rightarrow6$ \\
\code{smores\_physical\_driver\_to\_snake} & wheel/contact physics & 7 & 210 s / 0.002 s & four $6\rightarrow5\rightarrow6$ \\
\bottomrule
\end{tabularx}
\end{table}

\subsection{Robot Pack and authoring experiments}

Schema tests cover strict scalar types, ID grammar, unique catalogs, connector
orientation sets, policies, metadata, model-view recipes, and round-trip
serialization. Loader tests cover duplicate keys, version mismatch, split-document
issues, missing files, traversal, and symlinks. Writer tests inject failures during
staging and publication to check rollback and confirm that a reloaded pack is
semantically equal to the source.

URDF tests cover geometry, origins, axes, limits, multiple visuals/collisions, material
colors, package and local mesh resolution, draft asset rebasing, unsafe XML constructs,
and unresolved assets. Studio project tests verify functional mutation, referential
connector-type removal, connector parent-link validation, metadata persistence,
mapping cleanup, model-view CRUD, Save, reopen, and Export.

The committed \smoresep{} test loads the real asset tree, verifies zero validation
issues, confirms all four connector IDs and the corrected rear BOTTOM geometry, and
optionally compiles and docks two real modules through \mujoco{}.

\subsection{Docking semantic experiments}

The docking suite exercises mutual compatibility, all gender combinations, capture
shapes, antiparallel axes, discrete roll, tighter pair tolerances, connector-point
velocity, command policy, cooldown, loop closure, unsupported physical intent,
backend refusal, deterministic arbitration, release, overload, and assembly evolution.

Event expectations are exact rather than aggregate. A successful first connection
records candidate detection, dock commit, and assembly merge. A successful release
records undock commit and assembly split. Tests assert that a backend failure cannot
create a graph edge, and conformance tests compare committed decisions across installed
backends with gravity disabled.

\subsection{Two-module physical docking experiment}

The physical pair begins with two upright modules separated by a \SI{20}{mm} connector
gap. Gravity and the ground are enabled. After settling, the moving module drives its
PAN/TOP face into the target's rear BOTTOM face. The physical regression forbids any
post-creation root pose, twist, or wrench call and asserts actual wheel-joint travel.

\begin{table}[htbp]
\centering
\caption{Two-module physical docking acceptance criteria and outcomes.}
\label{tab:pair-results}
\begin{tabularx}{\textwidth}{L{0.38\textwidth}Y}
\toprule
Criterion & Enforced outcome \\
\midrule
Lifecycle & Exact dock, merge, undock, split semantic sequence \\
Topology & $0\rightarrow1\rightarrow0$ active edges; final two singleton assemblies \\
Failure counters & One dock, one undock, zero failures \\
Capture geometry & Connector translation at latch no greater than approximately \SI{1.001}{mm} \\
Release timing & Undock at least \SI{1.08}{s} after dock (hold plus release delay) \\
Actuation & Both moving wheel positions change by more than \SI{0.2}{rad} \\
Attitude & TILT below \SI{0.03}{rad}; PAN below \SI{0.01}{rad} in the checked final state \\
Separation & Mover reverses to root $x<-0.12$ m while retaining supported height \\
Physics ownership & Root pose/twist/wrench APIs fail the test if used after staging \\
\bottomrule
\end{tabularx}
\end{table}

The experiment establishes that the end-to-end actuation and connection path works for
one deterministic, pre-aligned approach. It does not establish arbitrary-pose docking
robustness.

\subsection{Kinematic Driver-to-Snake demonstration}

The kinematic seven-module demo isolates topology, events, graph behavior, and scenario
orchestration from contact. It stages the initial tree, executes each connector
replacement sequentially, and updates the Runtime Inspector. The topology begins with
seven nodes and six edges, alternates to five and back to six for each action, and ends
in the chain $1-3-2-4-5-6-7$. Gravity and ground are intentionally disabled.

This demonstration remains useful even after the physics scenario exists. It is faster,
deterministic, and can distinguish semantic regressions from control/contact regressions.

\subsection{Physical Driver-to-Snake experiment}

The physical experiment begins with the same six-edge tree staged upright at simulation
time zero. It then disables root motion injection through the regression guard and
executes all four replacements through effort-controlled wheel motion, ground contact,
release delays, measured capture, ideal welds, and paired contact exclusions.

\begin{table}[htbp]
\centering
\caption{Seven-module physical reconfiguration outcomes.}
\label{tab:snake-results}
\begin{tabularx}{\textwidth}{L{0.38\textwidth}Y}
\toprule
Measure & Result or enforced invariant \\
\midrule
Modules & 7 throughout \\
Connection history & $[6,5,6,5,6,5,6,5,6]$ \\
Final topology & Exact chain $1-3-2-4-5-6-7$ with six connector edges \\
Final assemblies & One assembly of size seven \\
Dock outcomes & 10 successful: six initial plus four replacements \\
Undock outcomes & 4 successful \\
Failures & 0 dock and 0 undock failures \\
Assembly events & 10 merges and 4 splits \\
Moving components & $\{1\}$, $\{7\}$, $\{1,2,3\}$, and $\{5,6,7\}$ \\
Component motion & Each wheel in the two three-module movers travels more than \SI{0.5}{rad} \\
Released-face clearance & Greater than \SI{20}{mm} for actions 3 and 4 \\
Completion & Test requires completion by 190 simulated seconds; reference trace about 177 s \\
Final hold & Stable and event-free through 210 simulated seconds \\
Wheel bound & No greater than $\pi/2+0.05\,\mathrm{rad/s}$ in the test \\
Allocation residual & Applied residual no greater than \SI{0.08}{m/s} \\
Pose health & Finite; unit quaternion; up projection $>\cos(0.15)$; root height 0.025--0.06 m \\
Physics ownership & Root pose/twist/wrench APIs forbidden after creation \\
\bottomrule
\end{tabularx}
\end{table}

The semantic event sequence is also exact: six initial
\code{DockCommitted}/\code{AssemblyMerged} pairs followed by four cycles of
\code{UndockCommitted}, \code{AssemblySplit}, \code{DockCommitted}, and
\code{AssemblyMerged}. The final hold verifies that completion does not leave stale
effort, accumulate new events, or destabilize the chain.

\subsection{Runtime visualization experiment}

Runtime Inspector tests use fake clocks and process hosts to check protocol framing,
initialization order, frame delivery, stderr handling, event continuity, graceful stop,
viewer-close behavior, and escalation. Presenter tests verify stable node layout,
parallel-edge curves, selection behavior, and stale-frame rejection. Qt tests run under
an offscreen/Xvfb display.

The current viewer-refresh regression executes ten \SI{2}{ms} physics steps while
synchronizing the native viewer twice, and a second test injects a \SI{30}{ms} sync cost
without causing an immediate repeated refresh. These are scheduling-correctness tests,
not graphics benchmarks.

\subsection{Playback speed}

At factor four, the documented approximately 177-s completion has a nominal wall-clock
target of $177/4\approx44.25$ s, and the 210-s display budget targets 52.5 s. These are
arithmetic targets, not measured achieved rates. The implementation never changes
solver step count: the full seven-module run remains 105,000 nominal \SI{2}{ms} steps.

\subsection{Result interpretation}

The experiments support three conclusions.

\begin{enumerate}
  \item A portable semantic package can drive both mock and physical-runtime workflows
        without moving connection meaning into backend code.
  \item Two-phase commit, canonical events, and derived graph views remain coherent
        through repeated physical dock and release operations.
  \item The backend-neutral joint command path is sufficient to build a nontrivial
        seven-module wheel/contact realization, including connected-component motion.
\end{enumerate}

They do not support claims of route robustness, calibrated hardware accuracy, large-N
scalability, or guaranteed four-times-real-time execution.

